% ==============================================================================
% core_step6_attitude.tex — SCI/SICE tutorial core deck, Step 6: attitude
% estimation (exercise 9). Code values from lesson_09_estimation/solution.cpp.
% SCI/SICE チュートリアル要点版 Step 6: 姿勢推定（実習 9）。
% ==============================================================================

% ------------------------------------------------------------------
\begin{frame}{ジャイロだけでは姿勢角が求まらない}
  \begin{columns}[c]
    \begin{column}{0.48\textwidth}
      \centering
      \adjustbox{max width=\columnwidth, max height=0.66\textheight}{\documentclass[border=10pt]{standalone}
\usepackage{tikz}
\usetikzlibrary{arrows.meta}

\begin{document}

\definecolor{truecolor}{RGB}{0, 120, 200}     % Blue - true value
\definecolor{driftcolor}{RGB}{220, 50, 30}     % Red - gyro integrated
\definecolor{gridcolor}{RGB}{210, 210, 210}
\definecolor{labelcolor}{RGB}{80, 80, 80}
\definecolor{annocolor}{RGB}{40, 160, 40}      % Green - annotations

\begin{tikzpicture}[
    axis/.style={-{Stealth[length=3mm]}, line width=1pt, labelcolor},
    trueline/.style={truecolor, line width=2.5pt},
    driftline/.style={driftcolor, line width=2pt},
    dim/.style={{Stealth[length=2.5mm]}-{Stealth[length=2.5mm]}, annocolor,
                line width=1pt},
    annolbl/.style={font=\normalsize\bfseries, annocolor},
  ]

  % Grid
  \foreach \x in {0, 1, ..., 8} {
    \draw[gridcolor] (\x, 0) -- (\x, 5);
  }
  \foreach \y in {0, 1, ..., 5} {
    \draw[gridcolor] (0, \y) -- (8, \y);
  }

  % Axes
  \draw[axis] (-0.3, 0) -- (8.8, 0) node[right, font=\large\bfseries] {$t$};
  \draw[axis] (0, -0.3) -- (0, 5.5) node[above, font=\large\bfseries] {$\theta$};

  % Time labels
  \foreach \x/\t in {0/0, 2/10, 4/20, 6/30, 8/40} {
    \node[font=\small, below, labelcolor] at (\x, -0.15) {\t s};
  }

  % True value (horizontal line at y=2)
  \draw[trueline] (0, 2) -- (8, 2);
  \node[font=\normalsize\bfseries, truecolor, below] at (3, 1.85) {True angle};

  % Gyro integrated value (drifting upward)
  \draw[driftline] plot[smooth, tension=0.5] coordinates {
    (0, 2.0) (0.5, 2.05) (1, 2.15) (1.5, 2.25) (2, 2.4)
    (2.5, 2.55) (3, 2.7) (3.5, 2.85) (4, 3.05)
    (4.5, 3.2) (5, 3.4) (5.5, 3.55) (6, 3.75)
    (6.5, 3.9) (7, 4.1) (7.5, 4.25) (8, 4.45)
  };
  \node[font=\normalsize\bfseries, driftcolor, above left] at (7, 4.3)
    {Gyro $\displaystyle\int\!\omega\,dt$};

  % Drift annotation at t=40s (x=8)
  \draw[annocolor, line width=0.8pt, dashed] (8, 2.0) -- (8, 4.45);
  \draw[dim] (8.4, 2.0) -- (8.4, 4.45);
  \node[annolbl, right] at (8.5, 3.2) {Drift};

\end{tikzpicture}
\end{document}
}
    \end{column}
    \begin{column}{0.50\textwidth}
      \footnotesize
      \begin{itemize}\setlength{\itemsep}{3pt}
        \item ジャイロは角速度 $\omega$ を測る。角度は積分 $\theta=\int\omega\,dt$ で求める
        \item 測定のバイアス（偏り）も積分されるので、時間とともに誤差が増える（ドリフト）
        \item 姿勢角そのものを目標にする制御（STABILIZE 以上）には、角度の推定が要る
      \end{itemize}
      \begin{alertblock}{そこで}
        \scriptsize 重力の向きを測る加速度計で角度を補う。加速度計は長期で安定だが、短期は振動や加減速でノイズが乗る
      \end{alertblock}
    \end{column}
  \end{columns}
\end{frame}

% ------------------------------------------------------------------
\begin{frame}{相補フィルタ: ジャイロと加速度計を混ぜる}
  \begin{columns}[c]
    \begin{column}{0.48\textwidth}
      \centering
      \adjustbox{max width=\columnwidth, max height=0.62\textheight}{\documentclass[border=10pt]{standalone}
\usepackage{tikz}
\usetikzlibrary{arrows.meta, positioning, calc}

\begin{document}

\definecolor{sfblue}{RGB}{0, 120, 200}
\definecolor{sfred}{RGB}{220, 50, 30}
\definecolor{sfgreen}{RGB}{40, 160, 40}
\definecolor{sfgray}{RGB}{80, 80, 80}
\definecolor{sflight}{RGB}{220, 240, 255}

\begin{tikzpicture}[
    node distance=1.5cm and 1.5cm,
    block/.style={draw, #1, line width=1.2pt, fill=#1!8,
                  minimum width=2.0cm, minimum height=0.9cm,
                  font=\normalsize\bfseries, rounded corners=3pt},
    block/.default=sfblue,
    gain/.style={draw, #1, line width=1.2pt, fill=#1!8,
                 circle, minimum size=0.9cm,
                 font=\normalsize\bfseries},
    gain/.default=sfblue,
    % Summing-junction circle diameter ~40% of block height (Block Diagram
    % Style convention; matches feedback_block.tex / pid_block.tex).
    % 加え合わせ点の円の直径はブロック高さの約4割（ブロック線図の流儀。
    % feedback_block.tex / pid_block.tex と揃える）。
    sum/.style={draw, sfblue, line width=1.2pt, fill=white,
                circle, minimum size=0.4cm,
                font=\normalsize\bfseries},
    input/.style={font=\normalsize\bfseries, sfgray},
    output/.style={font=\large\bfseries, sfblue},
    line/.style={-{Stealth[length=2.5mm]}, line width=1.2pt, sfgray},
    pathlbl/.style={font=\normalsize\bfseries},
  ]

  % --- Main path (gyro row): omega -> x(Delta t) -> Sigma1 -> alpha -> Sigma2 -> output ---
  % 主経路（ジャイロ行）: ω → ×Δt → Σ1 → α → Σ2 → 出力
  \node[input] (gyro_in) {Gyro $\omega$};
  \node[block=sfred, right=1.3cm of gyro_in] (dt_block) {$\times\Delta t$};
  \node[sum, right=1.3cm of dt_block] (sum1) {$\Sigma$};
  \node[gain=sfred, right=1.3cm of sum1] (alpha) {$\alpha$};
  \node[sum, right=1.3cm of alpha] (sum2) {$\Sigma$};
  \node[output, right=1.1cm of sum2] (output) {$\hat{\theta}$};

  % --- Accel path (top row): drops into Sigma2 from above, entirely outside
  %     the gyro feedback loop below (which stays self-contained around
  %     Sigma1). This is the layout fix versus the previous revision, where
  %     the accel row sat below the main row and visually nested inside the
  %     feedback loop even though it carries no feedback signal. ---
  % 加速度経路（上段）: Σ2 に上から入る。下側のジャイロ・フィードバック
  % ループ（Σ1周りで閉じる）とは完全に別の場所を通る。旧版は加速度経路が
  % 主経路の下段にあり，フィードバック信号ではないのに見た目上ループの
  % 内側に入り込んでいた。これを是正したのが今回のレイアウト。
  \node[input, above=2.0cm of gyro_in] (accel_in) {Accel};
  \node[block=sfgreen, right=1.3cm of accel_in] (atan) {atan2};
  \node[gain=sfgreen, right=1.3cm of atan] (one_alpha) {$1\!-\!\alpha$};

  % --- Feedback delay path (below row): previous estimate theta_hat_{k-1} ---
  % フィードバック遅延経路（下段）: 前回推定値 θ̂_{k-1} を戻す
  \node[block=sfblue, below=1.5cm of sum1, minimum width=1.6cm] (delay) {$z^{-1}$};

  % --- Branch point on the output line (theta_hat is tapped, not re-derived) ---
  % 出力線上の分岐点（θ̂ をそのまま分岐、再計算しない）
  \coordinate (tap) at ($(sum2)!0.5!(output)$);

  % --- Connections: gyro path ---
  \draw[line, sfred] (gyro_in) -- (dt_block);
  \draw[line, sfred] (dt_block) -- (sum1.west);
  \draw[line, sfred] (sum1.east) -- (alpha.west);
  \draw[line, sfred] (alpha.east) -- (sum2.west);

  % --- Connections: accel path. Drops straight down into Sigma2 from the
  %     north; a short right-angle jog (down, then across) absorbs the
  %     small misalignment between one_alpha's x-position and sum2's,
  %     since "right=1.3cm of" spacing is computed from each row's own
  %     node widths and the two rows use different-width nodes. ---
  % 加速度経路の接続。Σ2 に北側からまっすぐ下りる。one_alpha と sum2 の
  % x座標のわずかなズレ（"right=1.3cm of" の間隔は各行のノード幅から
  % 計算されるため，幅の違う行同士では厳密には揃わない）を、
  % 短い直角ジョグ（下りてから横に寄せる）で吸収する。
  \draw[line, sfgreen] (accel_in) -- (atan);
  \draw[line, sfgreen] (atan) -- (one_alpha);
  \draw[line, sfgreen] (one_alpha.south) -- ++(0,-0.4) -| (sum2.north);

  % --- Output line and its branch dot ---
  \draw[line, sfblue, line width=1.5pt] (sum2) -- (output);
  \fill[sfblue] (tap) circle (2.2pt);

  % --- Feedback path: branch dot -> straight down -> into delay from the
  %     east -> up into Sigma1 from below. With the accel path now routed
  %     above the main row, the lower half-plane belongs to this path
  %     alone, so (unlike the previous revision) no detour around another
  %     signal is needed. ---
  % フィードバック経路: 分岐点 → まっすぐ下へ → 遅延ブロックへ東側から
  % 入る → 上へ → Σ1 に南側から入る。加速度経路が上段に移ったことで
  % 下半分はこの経路専用になり，旧版で必要だった他信号を避ける迂回が
  % 不要になった。
  \draw[line, sfblue] (tap) -- (tap |- delay.east) -- (delay.east);
  \draw[line, sfblue] (delay.north) -- (sum1.south);

  % --- Sum node signs (Block Diagram Style convention): "+" sits above a
  %     horizontally (west-)entering arrow. For a vertically-entering arrow
  %     (into Sigma1 from below, into Sigma2 from above) the "+" sits to
  %     the arrow's right, mirroring Sigma1's existing bottom-entry sign. ---
  % 加え合わせ点の符号（ブロック線図の流儀）: 水平（西側）から入る矢印の
  % 「＋」はその上に置く。垂直に入る矢印（Σ1は下から，Σ2は上から）の
  % 「＋」はその右側に置く（Σ1の下からの入力の符号と同じ規則）。
  \node[font=\scriptsize, sfred,  anchor=south east] at ($(sum1.north)+(150:0.3)$) {$+$};
  \node[font=\scriptsize, sfblue, anchor=north west] at ($(sum1.south)+(330:0.3)$) {$+$};
  \node[font=\scriptsize, sfred,  anchor=south east] at ($(sum2.north)+(150:0.3)$) {$+$};
  \node[font=\scriptsize, sfgreen, anchor=south west] at ($(sum2.north)+(30:0.3)$) {$+$};

  % --- Path labels (kept off the signal lines) ---
  \node[pathlbl, sfred, below=0.2cm of dt_block] {High-pass};
  \node[pathlbl, sfgreen, above=0.2cm of atan] {Low-pass};
  \node[pathlbl, sfblue, below=0.15cm of delay] {前回値};

\end{tikzpicture}
\end{document}
}
    \end{column}
    \begin{column}{0.50\textwidth}
      \footnotesize
      \setlength{\abovedisplayskip}{2pt}\setlength{\belowdisplayskip}{2pt}%
      静止時の加速度計は重力の反力を測るので、傾きが求まる（$\phi$: ロール、$\theta$: ピッチ）:
      \[
        \begin{aligned}
        \phi_a&=\operatorname{atan2}(-a_y,-a_z)\\
        \theta_a&=\operatorname{atan2}\bigl(a_x,\sqrt{a_y^2+a_z^2}\bigr)
        \end{aligned}
      \]
      相補フィルタ（$\alpha=0.98$）:
      \[
        \hat\theta_k=\alpha(\hat\theta_{k-1}+\omega\,\Delta t)+(1-\alpha)\theta_a
      \]
      短期はジャイロ、長期は加速度計を使う（互いの弱点を補う）。ヨー角は重力から求まらない
    \end{column}
  \end{columns}
\end{frame}

% ------------------------------------------------------------------
\begin{frame}[fragile]{実習 9: 相補フィルタを実装し、機体の推定と比べる}
  \exgoal{自作の相補フィルタの角度を、機体既定の推定器 ESKF（誤差状態カルマンフィルタ）と比べる}
  \begin{columns}[T]
    \begin{column}{0.46\textwidth}
      {\footnotesize\textbf{\textcolor{sfblue}{手順（飛行しない）}}}\par\vspace{-2pt}
        \footnotesize
        \begin{enumerate}\setlength{\itemsep}{0pt}
          \item \code{sf lesson switch sci2026:9}
          \item \code{sf lesson edit} で右のコードにして保存
          \item \code{sf lesson build}、\code{sf lesson flash}
          \item \code{sf lesson monitor}、Teleplot で \code{cf\_roll} と \code{eskf\_roll} を重ねる
        \end{enumerate}
    \end{column}
    \begin{column}{0.52\textwidth}
      \begin{sflisting}[basicstyle=\ttfamily\scriptsize, moredelim={[is][\bfseries\color{sfred}]{<<}{>>}}]{user\_code.cpp（赤字を書く）}
<<float a_roll = atan2f(-ay, -az);>>
<<float a_pitch = atan2f(ax,>>
<<    sqrtf(ay*ay + az*az));>>
<<cf_roll = alpha*(cf_roll + gx*dt)>>
<<         + (1-alpha)*a_roll;>>
// pitch: same form (gy, a_pitch)
<<ws::print(">cf_roll:%.2f", ...);>>
<<ws::print(">eskf_roll:%.2f", ...);>>
      \end{sflisting}
    \end{column}
  \end{columns}
  \exresult{機体を手で傾けると 2 つの角度がほぼ一致して動く。$\alpha$ を変えて違いを見る}
  \exhelptodo{9}
\end{frame}
